\section{The Apollo Guidance Computer (AGC)}

The Apollo Guidance Computer (AGC) was a pioneering digital computer developed in the
early 1960s to navigate and control the Apollo spacecraft during their missions to the
Moon. Installed aboard both the Command Module and the Lunar Module, the AGC was
integral to the success of the Apollo program, providing real-time guidance,
navigation, and control capabilities.

Designed by the MIT Instrumentation Laboratory under the leadership of Charles Stark
Draper, the AGC featured a 16-bit word length, with 15 data bits and one parity bit,
and operated at a frequency of 2.048 MHz. The computer's memory was divided into 2,048
words of erasable magnetic-core memory (RAM) and 36,864 words of fixed core rope memory
(ROM), which stored the software and mission data. 

Astronauts interacted with the AGC through the Display and Keyboard (DSKY) interface, a
21-digit display and 19-button keyboard that allowed them to input commands and receive
information. This interface was crucial for manual adjustments and monitoring during
the missions. 

The AGC's software was written in AGC assembly language and stored on rope memory. The
bulk of the software was on read-only rope memory and thus could not be changed in
operation, but some key parts of the software were stored in standard read-write
magnetic-core memory and could be overwritten by the astronauts using the DSKY
interface, as was done on Apollo 14.


\subsection{Operating System}

The Apollo Guidance Computer (AGC) utilized a specialized operating system to manage
its resources and execute the software required for the Apollo missions. The operating
system, known as AGC software, was designed to perform real-time control and navigation
tasks with extreme reliability and efficiency. Unlike modern operating systems, which
handle general-purpose tasks for a wide variety of applications, the AGC's operating
system was dedicated exclusively to the guidance, navigation, and control of the Apollo
spacecraft.

One of the key innovations of the AGC operating system was its use of a
\textit{``priority scheduling"} system. This system ensured that higher-priority tasks,
such as navigation corrections or abort sequences, were processed before lower-priority
tasks, allowing the AGC to respond dynamically to changing conditions during the
mission. The operating system also utilized memory efficiently, with the available
2,048 words of RAM and 36,864 words of ROM being carefully allocated to critical
software routines.

The AGC operating system also incorporated error-checking mechanisms to maintain
reliability in the face of potential hardware malfunctions. Given the harsh and
unpredictable environment of space, the system was designed to operate autonomously,
with minimal input from the astronauts, while still allowing for manual intervention
when necessary. The operating system’s software was written in AGC assembly language
and was stored in the AGC’s rope memory, providing a highly optimized and
space-efficient solution for spacecraft operations.

This operating system consisted of two parts: the \indef{Executive}, a batch
job-scheduler that used priority based cooperative multi-tasking to manage longer
running jobs, and the \indef{Waitlist}, a timer interrupt-driven pre-emptive scheduler
that managed real-time tasks running on tight timing constraints. By design, waitlist
tasks were run in an interrupt-inhibited mode, ensuring uninterrupted processing. This
greatly limited how much computation these tasks are able to do, limiting each task to
about 150-200 instructions. As a result, a waitlist task requiring extensive
calculation would have to schedule other waitlist tasks, or a longer running executive
job. On the other hand, to meet the deadlines on critical tasks, executive jobs must
schedule waitlist tasks to be run within a certain deadline. 

The operating system also consisted of a robust recovery mechanism, that offered
different levels of restart protection for various tasks and jobs. \todo{elaborate}

The success of the AGC operating system relied on the cooperation between the executive
jobs and waitlist tasks and it's robust recovery system.


\section{The Executive}

The \indef{Executive} in the Apollo Guidance Computer is a multiprocessing preemptive
real-time operating system. It has a few features that set it apart from modern RTOS's,
and made it such a visionary piece of technology of its time. 

In this section, we'll give a brief overview of the features of the Executive.

\begin{itemize}
    \item Cooperative and preemptive
    \item Process construction:
        \item What we'd call process descriptors, were called Core Set. 
        \item Could ask for dyanmic storage at the point of process creation.
        \item Priority based scheduing
    \item Context switching: using the \texttt{NEWJOB} register
    \item Process sleep and wakeup
    \item Waitlist to do most of the real time operations

    \item Restart and recovery
\end{itemize}
